\section{Alpha particles in the coil field: do they stay in the plasma long enough?}\label{sec:alpha}

In the plasma section (Section~\ref{sec:plasma}) the alpha particles heat the plasma with \SI{611}{MW}. This assumes
that every alpha gives off its energy inside the plasma before it reaches the edge. Whether that holds depends on the
magnetic field. The user asked for a dynamic simulation with the magnetic field; this is the first part of it
(variant A): the orbits of 1000 alphas in the field of the project's own coils.

\paragraph{Terms used in this section.}
\begin{description}\itemsep2pt
\item[alpha particle] the helium nucleus created in every deuterium--tritium fusion, with \SI{3.5}{MeV} of kinetic
  energy. It is electrically charged and therefore follows the magnetic field lines.
\item[guiding centre] a charged particle circles rapidly around a field line (about \SI{4}{cm} radius for an alpha).
  Instead of computing every circle, one follows only the centre of this circle, the guiding centre. It moves along
  the field line and drifts slowly across it.
\item[passing and trapped] an alpha flying mainly along the field keeps going around the whole ring
  (\emph{passing}). One flying mainly across the field is reflected at places of stronger field, as by a mirror, and
  bounces back and forth (\emph{trapped}). Trapped particles are especially sensitive to small field errors.
\item[$v_\parallel/v$] the share of the velocity along the field: $+1$ or $-1$ means exactly along it, $0$ exactly
  across.
\item[slowing-down time] the time in which an alpha gives its energy to the plasma through collisions with electrons
  and ions, here about \SI{0.2}{s}. An alpha that stays in the plasma that long has done its job.
\item[last closed surface] the outermost of the nested flux surfaces, the edge of the confined plasma. An alpha that
  crosses it counts as lost here.
\item[quasisymmetry] a special form of stellarator field: although the machine is twisted, the field strength along
  the orbits looks as in a symmetric ring. Trapped particles then do not drift outwards on average. Our plasma shape is
  designed to be quasi-axisymmetric (QA).
\end{description}

\paragraph{Method.} The orbits are computed with simsopt \cite{simsopt} as collisionless guiding centres in the
vacuum field of the 16 coils from stage~2. Coils and shape are available at unit size ($R_0 = \SI{1.03}{m}$).
Collisionless orbits, however, depend only on the ratio of gyration radius to plasma radius ($\rho_L/a = 0.018$). So
instead of enlarging the coils, the velocity is reduced accordingly, and times and positions are scaled back to the
plant. The 1000 starting points follow the distribution of fusion power, and the directions are uniformly
distributed. The calculation runs to \SI{0.2}{s}, the slowing-down time. For the energy balance, each lost alpha is
assigned the energy it would still have at the time of loss according to classical slowing down \cite{stix1972}:
\[
u(t) = \bigl(u_0 + 1\bigr)\,e^{-3t/\tau_s} - 1,\qquad u = (E/E_c)^{3/2},\qquad \tau_s = \SI{0.195}{s},\ E_c = \SI{415}{keV}.
\]
The loss positions are assigned to a grid on the edge surface ($64\times32$ points, combined into $16\times8$ tiles of
about \SI{10}{m^2} each). This gives the heat flux.

\bild{alpha_verluste_en.png}{1.0}{Alpha particle losses: time history, who is lost, where they hit the edge, key
figures (1000 alphas, collisionless, field of the project's own coils).}
\lesehilfe{Top left: horizontally the time since the alphas were born, logarithmic from a hundredth of a millisecond
to \SI{0.2}{s}. Vertically, the percentage lost. Blue counts particles, orange their energy. The orange curve stays
lower because alphas lost late have already given off much of their energy. Hardly anything is lost in the first
millisecond; most losses happen between \SI{5}{ms} and \SI{100}{ms}. Top right: each dot is one alpha, horizontally its
direction at birth (0 = across the field), vertically its birth position (0 = plasma centre, 1 = edge). The red, lost
alphas lie almost all in a narrow strip around 0: they are trapped particles. Bottom left: the unrolled edge surface of
the plasma, horizontally once along the ring, vertically once around the cross-section. The colour is the heat flux,
the number the count of hits. The hits are bunched in one band, and twice over, because the machine repeats itself
twice. Bottom right: the key figures.}

\bild{alpha_bahnen_standbild_en.png}{1.0}{Three alpha orbits in the field of the project's own coils (still from the
video \texttt{docs/filme/alpha\_bahnen.mp4}): passing (blue), trapped (yellow), lost early (red).}
\lesehilfe{You are looking at the ring obliquely from above. The thin grey loops are cross-sections of the plasma edge
at 16 places; they show how the shape twists around the ring. Each coloured line is the path of an alpha's guiding
centre. The blue alpha flies along the field and keeps going around the whole ring; it stays safely inside. The yellow
one flies almost across the field and turns back at places of strong field; its path is bent back at the ends. The red
one is also a trapped alpha. Its orbit slowly moves outwards until it reaches the edge after \SI{0.13}{ms} (red dot
at the top). Orbits like this cause the losses in the figure above, though mostly much more slowly.}

\paragraph{Results.}
\begin{itemize}\itemsep1pt
\item \textbf{16.6\,\% of the alphas are lost} ($\pm1.2$\,\% statistical uncertainty with 1000 particles), only 10 of
  them in the first millisecond. Half of the losses have occurred after \SI{25}{ms}.
\item \textbf{They are almost all trapped alphas}: 95\,\% of the lost ones had $|v_\parallel/v| < 0.3$. On average they
  were born somewhat further out ($s = 0.24$ versus $0.18$).
\item \textbf{Lost energy: 11.6\,\% of the alpha power, i.e.\ \SI{71}{MW}}. The plasma misses this part of its
  self-heating; instead of \SI{611}{MW} about \SI{540}{MW} remain.
\item \textbf{The losses arrive bunched.} Averaged over the whole edge it is only \SI{58}{kW/m^2}, but almost all hits
  lie in one band in the upper half of the cross-section, in two spots. In the hottest tile of about \SI{10}{m^2} it
  is around \SI{2.4}{MW/m^2}. The first wall of this project is designed for \SI{0.5}{MW/m^2} (cooling section).
  Because of the coarse tiles and only 27 hits in the hottest tile, the peak value is uncertain; the true peak is more
  likely higher.
\end{itemize}

\paragraph{Assessment.} The losses do not come as a prompt loss in the first transits but slowly, over milliseconds up
to the slowing-down time. This is the typical picture of trapped particles drifting outwards in a field that is not
quite quasisymmetric. The plasma shape itself is very precisely quasisymmetric (quasisymmetry objective in DESC
$2.4\cdot10^{-5}$, smaller than the reference \emph{precise QA} with $5.8\cdot10^{-5}$). The coils, however, only
approximate this shape: the field is inclined to the edge surface by 0.27\,\% on average and 0.96\,\% at most
($|B_n|/|B|$). Whether most of the losses come from these coil errors or are already inherent in the shape has not yet
been checked. Landreman and Paul showed that very precise quasisymmetry is achievable and then gives excellent
confinement \cite{landreman2022}. A loss of more than 10\,\% of the alpha energy would be too much for a power plant,
above all because of the bunched wall load.

\paragraph{Limits.}
\begin{itemize}\itemsep1pt
\item \textbf{Collisionless}: in reality the orbits change while the alpha slows down and changes direction through
  collisions. Collisions can turn trapped particles into passing ones and vice versa; this can shift the losses either
  way.
\item \textbf{Vacuum field}: currents in the plasma (finite pressure) are not included.
\item \textbf{The loss location is the plasma edge}, not the wall. The wall lies further out, and the hits spread
  somewhat differently there.
\item \textbf{Statistics}: 1000 alphas are enough for the loss fraction, but not for a fine wall-load map.
\end{itemize}

\paragraph{Next steps.} (1) Compute the same alphas in the field of the plasma shape alone (without coil errors) to
separate the causes. (2) If the coils are the cause, re-optimise the coils with a target for alpha losses. (3) More
particles and collisions for a reliable wall-load map.
